\section{KDE Project Analysis} % {{{

\marginnote{\centering\includegraphics[width=0.5\marginparwidth]{kde/logo}}

KDE is a desktop environment designed to run on UNIX based systems as well as
Microsoft Windows and Mac OS X systems \cite{KDEPress,KDEAbout}. Since 2010 and
the release of KDE 4.4 the project is known as \ac{KDE SC}. It is composed by
the desktop environment named Plasma Desktop and several core application for
the daily needs. \ac{KDE SC} is licensed under the \ac{GPL} and \ac{LGPL}. The
name was originally an acronym for Kool Desktop Environment and later for K
Desktop Environment, however it is no longer in use. KDE is a modular project
consisting of several libraries which allow developer to build their programs
around the platform. As the KDE project and the \ac{KDE SC} stand for a wide
range of technologies and applications, this analysis will be limited to the
core of the \ac{KDE SC}, also known as KDE Base.

\subsection{History} % {{{

KDE was first announced started by Matthias Ettrich in 1996 as a desktop
environment for end users with the same feel and look for all applications
\cite{KDEAnnouncement}. The name KDE was more a word play with the existing
Common Desktop Environment (CDE), which was highly popular at that time. As the
graphical library, Matthias Ettrich chose the Qt framework which was and is
developed by the company Trolltech. KDE 1.0 was then released in 1998
\cite{KDEHistory}. At the same time, Trolltech dual licensed the framework
under the Q Public License (QPL) and a proprietary software license. However,
it was debated if that license was compatible with the \ac{GPL} and in 2000 the
framework was licensed under the \ac{GPL} which ceased the criticism.

In 2009 the KDE project renamed the project to \ac{KDE SC} and redefined the
project as a community which delivers free software for user interfaces by
emphasizing the KDE technologies \cite{KDESC}. All software created with KDE
technologies will be a KDE project. However \ac{KDE SC} only contains projects
which come out from the project itself and share a common release cycle.

\ac{KDE SC} 4 was a new approach to the desktop metaphor and created a new user
experience for users. The centerpiece is the Plasma Workspace which exists for
several devices, such as for desktop computers, netbooks, tablets and
smartphones.

% }}}

\subsection{Community} % {{{

The KDE project is one of the largest FOSS projects and according to the
project the second largest after the Linux Kernel \cite{KDEPress}. There are
about 1800 active contributors to the project and it's surroundings and is used
by over a million people. Also, the community spans around projects basing on
KDE technologies.

Most of the communication in the KDE project is done via mailing lists, most
importantly \emph{kde-devel} and \emph{kde-core-devel}
\cite{KDEProjectManagement,KDEContribute}. While the first one is mostly used
for communication by application developers, the latter one is for
communication of the KDE core project. Furthermore communication can happen
over \ac{IRC} channels, blogs or forums.

The most important KDE conference is the annual \emph{Akademy} which is hold at
different locations in Europe \cite{KDEHistory}. The first KDE conferences were
however named after the current release which started with \emph{KDE One} in
1997. That is also the reason, why the conferences did not happen annually at
that time. Not until 2004 it changed the name to Akademy. Beginning with 2009
the project incorporates the conference into the \emph{Desktop Summit} which is
a collaboration with the GNOME project.

\begin{table}
  \centering
  \begin{tabularx}{\textwidth}{lXl}
    \toprule
    \tableheadline{Event}   & \tableheadline{Venue}           & \tableheadline{Date} \\
    \midrule
    KDE One                 & Arnsberg, Germany               & 1997 \\
    KDE Two                 & Erlangen, Germany               & 1999 \\
    KDE Three Beta          & Trysil, Norway                  & 2000 \\
    KDE Three               & Nürnberg, Germany               & 2002 \\
    Kastle                  & Nové Hrady, Czech Republic      & 2003 \\
    aKademy                 & Ludwigsburg, Germany            & 2004 \\
    aKademy                 & Málaga, Spain                   & 2005 \\
    aKademy                 & Dublin, Ireland                 & 2006 \\
    aKademy                 & Glasgow, Scotland               & 2007 \\
    Akademy                 & Sint-Katelijne-Waver, Belgium   & 2008 \\
    Desktop Summit          & Gran Canaria, Spain             & 2009 \\
    Akademy                 & Tampere, Finland                & 2010 \\
    Desktop Summit          & Berlin, Germany                 & 2011 \\
    Akademy                 & Tallinn, Estonia                & 2012 \\
    \bottomrule
  \end{tabularx}
  \caption{Previous and planned KDE conferences}
\end{table}

Owing to the size of the KDE project it is organized around many independent
teams with a communication structure between them. There are few groups who
stand above those teams and only take care about coordinating the project
\cite{KDEDevelopmentModel,KDEProjectManagement}.

\begin{description}

  \item[Contributor] After having contributed to any KDE project for some time
    and planning to contribute in the future, a KDE Contributor account can be
    requested \cite{KDEContribute,KDEContributor}. In the request one must
    state why they are interested in having an account. The KDE sysadmin team
    will then check the application and grant an account or not.

  \item[Module Coordinator] As the KDE project is highly modular, quite each
    application or library has an independent team with it's own structure. A
    module coordinator plans and executes releases and the direction of a
    single module. As the teams are so diverse, there are many ways to become a
    module coordinator. The most obvious is of course by unweary contributions
    to a project.

\begin{figure}[htbp]
  \centering
  \includegraphics[width=0.8\textwidth]{kde/commits_by_author}
  \caption{Monthly activity of the most active KDE core developers}
\end{figure}

  \item[Core Team] One of the most important teams is the Core Team as it
    defines the overall direction the project is heading
    \cite{KDEProjectManagement}. There is no single person inside the team
    which takes the decision, instead the team takes a decision by discussing
    the issue and coming up with a solution. The primary communication method
    is the kde-core-devel mailing list which is publicly readable, however an
    approval is required to join it. The KDE project does not define clear
    restrictions to membership and so a membership is by invite only after
    distinguishing or outstanding work for the KDE project.

  \item[Release Team] The actual release and the release schedules are provided
    and enforced by the Release Team \cite{KDEReleaseTeam}. It is composed by
    module coordinators and other release team members who actually do the
    releases. The Release Team makes sure that all modules are following the
    release schedule and are in a good shape. Also they decide on code freeze
    breaks and take decision about future features of a release.

  \item[KDE Founder] Matthias Ettrich, the founder of the KDE project is no
    longer actively involved in the project as he works full time on the
    underlying Qt framework. However he has still an important voice in the
    community and contributes indirectly to KDE by his role in the Qt project.

\end{description}

\begin{figure}[htbp]
  \centering
  \includegraphics[width=0.8\textwidth]{kde/commits_by_year}
  \caption{Yearly overview of commits to KDE base}
\end{figure}

% }}}

\subsection{Release Process} % {{{

The KDE project uses a version naming scheme with three numbers and
distinguishes between major and minor releases
\cite{KDEReleaseTeam,KDEReleaseSchedule,KDESchedule}. The numbers are called
major, minor and micro numbers. The major and minor number define major
releases while the micro number defines a maintenance release. Depending on
whether the major or minor number changes the release is a platform or standard
release. A platform release defines a new series of releases which can break
backwards compatibility to the previous platform release. They are often
planned a long time beforehand with usually also changing the used libraries or
switching their used versions. A standard release however has to maintain
backwards compatibility with it's series. They can have new features and user
interface changes, however a KDE application has to work on all releases of the
same platform. Maintenance releases are not allowed to have new features,
however bug fixes and small enhancements are allowed. Beginning with \ac{KDE
SC} 4 the cycle changed to a major release every six months and a maintenance
release roughly every month except when a major release comes out.

\begin{figure}[htbp]
  \centering
  \includegraphics[width=\textwidth]{kde/releases}
  \caption{Platform and Standard Releases of KDE}
\end{figure}

\begin{figure}[htbp]
  \centering
  \includegraphics[width=\textwidth]{kde/punchcard}
  \caption{Time based view on commits of Core contributors}
\end{figure}

To ensure the release schedule, the release team publishes a new release
schedule after each major release. The schedule includes all deadlines and
freezes which will lead to the next major release and includes future
maintenance releases. This means of course that two release schedules will
overlap until the month before the new major release. The following freezes and
releases are provided and enforced by the release team
\cite{KDEReleaseSchedule}.

\begin{description}

  \item[Soft Feature Freeze] Approximately three weeks before the first beta
    release no new features are allowed if they are not already approved as
    planned features. Not ready features have to be postponed to the next major
    release.

  \item[Dependency Freeze] Approximately two weeks before the first beta
    release no new or new dependency versions are allowed anymore. However it
    is possible to request an exception by the release team.

  \item[Soft Message, Soft API and Hard Feature Freeze] One week before
    the first beta release no new strings and changes to the existing messages
    can be made except corrections. Additionally the existing API should be
    almost finished and if changes are made they should be reported to the
    relevant teams. Lastly no new features can be added to the projects, even
    if they were planned for this release cycle.

  \item[Beta Release] Approximately two months before the major release the
    first beta release gets published. This is usually followed by a second
    beta release two weeks later.

  \item[Hard API, Message and Documentation Freeze] Five weeks before the
    major release the \ac{API}, translatable messages and the documentation
    must be finished and can't be changed anymore.

  \item[Release Candidate] Approximately a month before the final release the
    first release candidate will be published. Two weeks later the second
    release candidate follows.

  \item[Final Release] The final release will then be published three weeks after
    the last release candidate.

  \item[Minor Releases] A new minor release will be published each following
    month until the next major release gets published. This normally leads to
    four minor releases.

\end{description}

\begin{figure}[htbp]
  \centering
  \includegraphics[width=0.8\textwidth]{kde/commits_by_month}
  \caption{Amount of commits per month of Core contributors}
\end{figure}

% }}}

\subsection{Development} % {{{

During the development of a new major release, new project goals are set using
feature proposals \cite{KDEDevelopmentModel,KDEFAQ}. This normally happens
after each major release when KDE developers have the time to propose new
features for the next major release. Each feature must have an author who is
willing to do the implementation of that feature. This often limits the
features to projects of already contributing KDE developers. A KDE developer is
free to implement the feature by himself, however it often is accompanied by a
discussion with the specific team and several adjustments to the
implementation.

At the time of the soft feature freeze it is decided whether the feature will
be available for the upcoming major release or if it will be postponed to the
next major release. If the implementation already started and it is almost
finished, it can continue. Otherways it will be postponed. This development
process is a quite deliberate approach where people are free to implement what
they see missing in the project.

\begin{figure}[htbp]
  \centering
  \includegraphics[width=0.8\textwidth]{kde/authors_by_month}
  \caption{Amount of distinct authors of KDE base over time}
\end{figure}

% }}}

% }}}
